The contribution of the present paper is explicit and layered.  We do not
claim the stability-enhanced rank--trace inequality or the seven-point
architecture as new.  We use the already certified nonuniform pressure vector
\begin{equation}\label{eq:p}
 p=\frac1{10^7}(2714,3733,3553,3553,3733,2714),
 \qquad
 \beta:=\sum_{j=1}^6p_j=\frac1{500},
\end{equation}
for which the seven-point functional satisfies
\begin{equation}\label{eq:local-cert-intro}
 \mathcal F_p(g_1,\ldots,g_6)\ge\frac{39}{10000}
 \qquad(g_j\ge0).
\end{equation}

The new argument has two stages.  First, rather than replacing the accumulated
position-dependent pressure inside a block by
\(\beta\,\operatorname{span}(B)\), we retain its exact coefficients until the
final average over shifted blocks.  This reduces the global pressure mass from
the coarse \(\beta(m-1)\) to the exact \(\beta(m-6)\).  Already at \(m=262\)
this gives
\[
 \frac{655000H_{\rm MT}-1280}{652504}
 =
 0.6731115225500757511923586\ldots .
\]
Second, an adjacent-pair pinching argument supplies an independent lower bound
for the Gram defect in terms of the consecutive-gap kernel energies.  Together
with one rigorously enclosed scalar value of the Montgomery--Taylor kernel,
this permits \(m=450\) and yields Theorem~\ref{thm:main}.  No new
seven-point local certificate is required.
